\section{Introduction}
\label{sec:introduction}
Speculative decoding uses a draft to propose tokens and a target to evaluate several positions at once~\cite{leviathan2023}. On consumer systems, this extra work within a request can amortize target execution when memory capacity and request traffic limit batching across requests. The implementation also changes how the target executes: scalar continuation becomes multi-position verification, cache contents follow a different history, and rejected positions require reconciliation. A valid sampling rule therefore leaves two practical questions open: does the quantized runtime reproduce the required output, and does it reduce the cost of producing it?

We study systems in which ordinary greedy decoding under a fixed runtime configuration defines the reference. Such a requirement arises in regression tests and applications that need repeatable outputs. Reference preservation means matching the registered token sequence or endpoint; it does not establish that the answer is functionally correct or that every intermediate floating-point value is identical. The runtime designer must choose where to enforce consistency and then measure the resulting cost. A constraint can repair one input while breaking another, or recover the output at the expense of the intended acceleration.

The visible token prefix does not determine the necessary constraint. Two executions can emit identical tokens while retaining different key/value (KV) values, pending tokens or sampler history. Their next decision can follow the retained state, the current execution recipe, or an interaction between the two. Conversely, different scores can produce the same immediate greedy decision while later outputs diverge. Distinguishing these cases helps decide whether to constrain state formation, current verification, or both.

A repair may improve the target path before a draft is added. Comparing repaired speculation only with ordinary decoding measures the combined deployment change. A matched target-only execution, using the repaired path without a draft, measures how much additional benefit speculation supplies. Both paths must be executed and charged under the same output requirement; acceptance rates and target-call counts cannot replace their measured costs.

An RTX continuation intervention repairs one task and regresses another, leaving aggregate reference agreement unchanged. A later phase-scoped attention implementation qualifies full responses on 24 inputs unused in the accessible project history. Its resident-request speed is $1.742\times$ ordinary decoding, but the matched draft-free path already achieves $1.683\times$; adding speculation gives $1.035\times$. A distinct B570 implementation preserves its registered endpoint while remaining substantially slower than both draft-free baselines. We organize the study around three questions:
\begin{enumerate}
\item \RQtag{1}, diagnosis: does the observed deviation follow retained state, the current recipe, or their interaction?
\item \RQtag{2}, repair: which phase constraints preserve the reference within a defined implementation family?
\item \RQtag{3}, cost: after qualification, what benefit remains against ordinary decoding and the matched target path?
\end{enumerate}

The contribution is a linked diagnosis, repair and cost evaluation within explicit output contracts. Natural donor-state crosses with own-path controls distinguish local state, path and interaction effects on CUDA and Vulkan; a separate same-state B570 fork study establishes the limits of immediate agreement. The RTX query-layout transformation retains batched projection and feed-forward work, and a complete four-switch factorial identifies a sufficient target-phase scope within the tested family. Measured ordinary/matched/speculative comparisons separate the full deployment gain from the speculative increment. B570 supplies a qualified-but-slower counterexample, while the auxiliary M4 and 358H cases delimit the value of added control and draft-only optimization.

The study concerns runtime design and finite qualification. Its crosses are offline experiments, and the repair remains specific to the tested Qwen3 family and configurations. Sections~\ref{sec:runtime} and~\ref{sec:methods} define the implementation context and measurement contracts; Sections~\ref{sec:diagnosis}--\ref{sec:cost} answer the three questions. Section~\ref{sec:implications} draws the resulting design choices and states their limits. The supplement retains campaign details, historical results and the complete evidence trail.
